\section{Curriculum Package --- Python, Java, and
JavaScript}\label{curriculum-package-python-java-and-javascript}

This package contains three complete introductory programming curricula,
each adapted for visually impaired students who use screen readers.

\begin{center}\rule{0.5\linewidth}{0.5pt}\end{center}

\section{Package Contents}\label{package-contents}

\begin{verbatim}
curriculum/
|--- README.md                          (this file)
|--- python/
|   |--- unit00-setup.md
|   |--- unit01-foundations.md
|   |--- unit02-numeric-computations.md
|   |--- unit03-control-flow.md
|   |--- unit04-06-data-functions.md
|   |--- unit07-09-oop-files.md
|   +--- unit10-11-projects-capstone.md
|--- java/
|   |--- README.md
|   |--- unit00-01-setup-foundations.md
|   |--- unit02-04-io-control-arrays.md
|   |--- unit05-08-strings-methods-oop-files.md
|   +--- unit09-11-collections-projects.md
+--- javascript/
    |--- README.md
    |--- unit00-01-setup-foundations.md
    |--- unit02-04-io-control-data.md
    |--- unit05-08-strings-functions-oop-files.md
    +--- unit09-11-modules-collections-projects.md
\end{verbatim}

\begin{center}\rule{0.5\linewidth}{0.5pt}\end{center}

\section{Curriculum Comparison}\label{curriculum-comparison}

{\def\LTcaptype{none} % do not increment counter
\begin{longtable}[]{@{}
  >{\raggedright\arraybackslash}p{(\linewidth - 6\tabcolsep) * \real{0.1604}}
  >{\raggedright\arraybackslash}p{(\linewidth - 6\tabcolsep) * \real{0.3019}}
  >{\raggedright\arraybackslash}p{(\linewidth - 6\tabcolsep) * \real{0.2830}}
  >{\raggedright\arraybackslash}p{(\linewidth - 6\tabcolsep) * \real{0.2547}}@{}}
\toprule\noalign{}
\begin{minipage}[b]{\linewidth}\raggedright
Feature
\end{minipage} & \begin{minipage}[b]{\linewidth}\raggedright
Python
\end{minipage} & \begin{minipage}[b]{\linewidth}\raggedright
Java
\end{minipage} & \begin{minipage}[b]{\linewidth}\raggedright
JavaScript
\end{minipage} \\
\midrule\noalign{}
\endhead
\bottomrule\noalign{}
\endlastfoot
\textbf{Difficulty} & Beginner & Intermediate & Beginner-Intermediate \\
\textbf{Typing} & Dynamic & Static & Dynamic \\
\textbf{Syntax} & Minimal & Verbose & Moderate \\
\textbf{Best For} & First language, data science, AI & Enterprise,
Android, CS theory & Web development, full-stack \\
\textbf{Terminal Tool} & \texttt{python3\ file.py} & \texttt{javac} +
\texttt{java} & \texttt{node\ file.js} \\
\textbf{Testing} & \texttt{unittest} & JUnit 5 & Jest \\
\textbf{JSON} & \texttt{json} module & Gson library & Built-in
\texttt{JSON} object \\
\textbf{Input} & \texttt{input()} & \texttt{Scanner} & \texttt{readline}
module \\
\end{longtable}
}

\begin{center}\rule{0.5\linewidth}{0.5pt}\end{center}

\section{Accessibility Notes (All
Languages)}\label{accessibility-notes-all-languages}

All three curricula share these accessibility principles:

\begin{enumerate}
\def\labelenumi{\arabic{enumi}.}
\tightlist
\item
  \textbf{Screen Reader Optimized Mode in VS Code:} Press
  \texttt{Ctrl\ +\ Shift\ +\ P} -\textgreater{} ``Toggle Screen Reader
  Optimized''.
\item
  \textbf{Indentation Announcement:} Configure your screen reader to
  announce indentation changes --- critical for Python's block
  structure; helpful for Java and JavaScript too.
\item
  \textbf{Error Navigation:} Use \texttt{F8} / \texttt{Shift+F8} to jump
  between errors without scanning visually.
\item
  \textbf{All Output is Text-Based:} No graphical output required for
  any core exercises.
\item
  \textbf{Terminal First:} All programs run from the integrated
  terminal, which is accessible with screen readers.
\end{enumerate}

\begin{center}\rule{0.5\linewidth}{0.5pt}\end{center}

\section{Recommended Learning Paths}\label{recommended-learning-paths}

\subsection{Complete Beginner}\label{complete-beginner}

Start with \textbf{Python} -\textgreater{} then \textbf{JavaScript}
(many concepts transfer directly).

\subsection{Web Developer Track}\label{web-developer-track}

Start with \textbf{JavaScript} -\textgreater{} deepen with Node.js and
frameworks.

\subsection{Computer Science Theory
Track}\label{computer-science-theory-track}

Start with \textbf{Python} -\textgreater{} move to \textbf{Java} for
stronger OOP and typing foundations.

\subsection{All-Rounder}\label{all-rounder}

\textbf{Python} (learn concepts) -\textgreater{} \textbf{Java} (learn
strict typing and OOP) -\textgreater{} \textbf{JavaScript} (apply to the
web).

\begin{center}\rule{0.5\linewidth}{0.5pt}\end{center}

\section{Credits}\label{credits}

Curriculum designed for accessibility-first computer science
education.\\
Adapted from:

\begin{itemize}
\tightlist
\item
  Matthes, E. (2023). \emph{Python Crash Course} (3rd ed.). No Starch
  Press.
\item
  Zelle, J.M. (2017). \emph{Python Programming: An Introduction to
  Computer Science} (3rd ed.). Franklin, Beedle \& Associates.
\item
  Eck, D.J. (2022). \emph{Introduction to Programming Using Java} (9th
  ed.). {[}Free at math.hws.edu/javanotes{]}
\item
  MDN Web Docs. {[}developer.mozilla.org{]}
\end{itemize}

Original Python curriculum by Michael Ryan Hunsaker, M.Ed., Ph.D.~---
Davis School District.
